\documentclass[11pt]{article}
\usepackage[utf8]{inputenc}
\usepackage{amsmath,amssymb,amsthm}
\usepackage[margin=1in]{geometry}
\usepackage{microtype}
\usepackage{hyperref}
\usepackage{xcolor}
\usepackage{enumitem}
\usepackage{newunicodechar}
\emergencystretch=3em
\newunicodechar{ℝ}{\ensuremath{\mathbb{R}}}
\newunicodechar{ℕ}{\ensuremath{\mathbb{N}}}
\newunicodechar{∫}{\ensuremath{\int}}
\newunicodechar{λ}{\ensuremath{\lambda}}
\newunicodechar{α}{\ensuremath{\alpha}}
\newunicodechar{γ}{\ensuremath{\gamma}}
\newunicodechar{∂}{\ensuremath{\partial}}
\newunicodechar{²}{\ensuremath{^2}}
\newunicodechar{³}{\ensuremath{^3}}
\newunicodechar{·}{\ensuremath{\cdot}}
\newunicodechar{≤}{\ensuremath{\leq}}

\newcommand{\tideref}[2][]{\href{https://therisensea.org/tides/#2/}{\ifx&#1&\texttt{#2}\else#1\fi}}
\newcommand{\experimentref}[2][]{\href{https://therisensea.org/experiments/#2/}{\ifx&#1&\texttt{#2}\else#1\fi}}

\title{Tide retrospective:\\\emph{higher \(h\)-derivatives of the anharmonic partition function}}
\author{Claude Opus 4.8 (1M ctx), with Daniel Murfet}
\date{2026-05-30}

\begin{document}
\maketitle

\begin{abstract}
For the 1D anharmonic potential $L(x) = \tfrac{\lambda}{2}x^2 +
\tfrac{\alpha}{6}x^3 + \tfrac{\gamma}{24}x^4$ ($\alpha^2 < 3\lambda\gamma$)
under the linear perturbation $A(x)=x$, the perturbed partition function
is $Z(h) = \int e^{-t(L(x)+hx)}\,dx$. The existing regularity instance
proved its first $h$-derivative at $h=0$. This tide proves the general
iterated-derivative \emph{step}: for every $n\in\mathbb{N}$,
\[
  \frac{d}{dh}\Big|_{h=0}\!\int (-(tx))^n e^{-t(L(x)+hx)}\,dx
    = \int (-(tx))^{n+1} e^{-t L(x)}\,dx,
\]
whose $n=1,2$ instances are the second and third $h$-derivatives of $Z$ at
$h=0$ ($\int (tx)^2 e^{-tL}$ and $\int -(tx)^3 e^{-tL}$). New file
\texttt{AnharmonicPartitionHigherDeriv.lean} ($\approx$200 lines), commit on
branch \texttt{tide/anharmonic-partition-higher-deriv}; builds clean, no
\texttt{sorry}/\texttt{axiom}/\texttt{native\_decide}, no lint warnings;
both derivative forms numerically confirmed to $\sim$$10^{-6}$. Directly
reuses the general $|x|^m$ integrability witness and the two public pointwise
helpers landed in the FDT-capstone tide earlier the same day.
\end{abstract}

\section{Setting}

The laplace seabed formalises the Laplace/Gibbs calculations of the SLT
Susceptibility Primer. A recent tide closed the last \texttt{sorry} in the
anharmonic \texttt{GibbsRegularity} instance, whose third field,
\texttt{partition\_hasDerivAt}, gives the \emph{first} $h$-derivative at
$h=0$ of the perturbed partition $Z(h) = \int e^{-t(L(x)+hx)}\,dx$, namely
$\int (-tx)\,e^{-tL}$, proved by Mathlib's dominated-differentiation lemma
\texttt{hasDerivAt\_integral\_of\_dominated\_loc\_of\_deriv\_le}. Earlier
this same day the FDT-capstone tide\footnote{\tideref[the anharmonic FDT
capstone]{2026-05-30-tide-anharmonic-fdt-capstone}.} added a general
integrability witness $\int |x|^m e^{-tL}<\infty$ for all $m$ and made two
pointwise helpers public.

The minimal good target outward from there: the higher $h$-derivatives. The
flow-equation / fourth-cumulant identity needs the second and third
derivatives of $Z$ at the unperturbed point $h=0$. Rather than prove $n=2$
and $n=3$ as two separate lemmas, a single general-$n$ step lemma covers
both (and all higher orders) at no extra cost, because the new $|x|^m$
witness is already uniform in the power.

\section{The thing formalised}

\begin{quote}\itshape
For $\lambda,\gamma>0$, $\alpha^2<3\lambda\gamma$, $t>0$, and every
$n\in\mathbb{N}$,
\[
  \frac{d}{dh}\Big|_{h=0}\!\int_{\mathbb{R}} (-(tx))^n\,
      e^{-t(L(x)+hx)}\,dx
    \;=\; \int_{\mathbb{R}} (-(tx))^{n+1}\,e^{-t L(x)}\,dx .
\]
\end{quote}

The Lean signature:

\begin{verbatim}
theorem anharmonic_partition_deriv_step
    {lam alpha gamma t : ℝ} (n : ℕ)
    (hlam : 0 < lam) (hgamma : 0 < gamma)
    (hdisc : alpha ^ 2 < 3 * lam * gamma) (ht : 0 < t) :
    HasDerivAt
      (fun h => ∫ x, (-(t * x)) ^ n *
        Real.exp (-(t * (anharmonicPotential lam alpha gamma x + h * x))))
      (∫ x, (-(t * x)) ^ (n + 1) *
        Real.exp (-(t * anharmonicPotential lam alpha gamma x)))
      0
\end{verbatim}

The two named corollaries \texttt{anharmonic\_partition\_secondDeriv}
($n=1$, right-hand side $\int (tx)^2 e^{-tL}$) and
\texttt{anharmonic\_partition\_thirdDeriv} ($n=2$, right-hand side
$\int -(tx)^3 e^{-tL}$) are the literal second and third derivatives the
direction asked for. They are the function being differentiated read as the
$(n)$-th derivative integrand: differentiating $\int(-(tx))^1 e^{-t(L+hx)}$
at $0$ yields $Z''(0)$, and differentiating $\int(-(tx))^2 e^{-t(L+hx)}$
yields $Z'''(0)$.

\section{Proof strategy}

The proof is the $n=1$ \texttt{partition\_hasDerivAt} field with one extra
constant-in-$h$ factor $(-(tx))^n$ carried through. Feed
\texttt{hasDerivAt\_integral\_of\_dominated\_loc\_of\_deriv\_le} the family
$F(h,x)=(-(tx))^n e^{-t(L(x)+hx)}$ on the unit ball around $0$:

\begin{itemize}[itemsep=2pt]
\item \textbf{Pointwise derivative.} The public
  \texttt{anharmonic\_perturbed\_pointwise\_hasDerivAt} differentiates the
  Boltzmann factor; \texttt{HasDerivAt.const\_mul $(-(tx))^n$} multiplies in
  the monomial, and \texttt{ring} (with the exponential as an opaque atom)
  collapses $(-(tx))^n\cdot(-(tx))$ to $(-(tx))^{n+1}$ via \texttt{pow\_succ}.
\item \textbf{Dominator.} Writing $F'(h,x)=(-(tx))^n\cdot[(-tx)\,e^{-t(L+hx)}]$
  and reusing the $n=0$ pointwise bound
  \texttt{anharmonic\_perturbed\_pointwise\_bound} on the bracket, the
  dominator is $t^{n+1} e^{t/2c} |x|^{n+1} e^{-(t/2)L}$, integrable by the
  general $|x|^{n+1}$ witness at half temperature.
\item \textbf{Endgame.} The lemma returns the derivative as
  $\int F'(0,x)$, whose integrand still carries a $+0\cdot x$; an
  \texttt{integral\_congr\_ae} with \texttt{ring} under the integral drops it.
\end{itemize}

\section{Roadblocks and resolutions}

One real bug, caught by the build, then by re-reading the magnitude. The
first draft tried to dominate $F(0,x)$ — needed for the \emph{integrability
of $F$ at $0$} hypothesis — by the bare witness $|x|^n e^{-tL}$. But
$\|(-(tx))^n e^{-tL}\| = t^n |x|^n e^{-tL}$ carries a $t^n$ the bare witness
lacks, so the domination inequality was simply false for $t>1$ (and the
\texttt{rw} chain that was supposed to massage it failed to match, which is
how Lean flagged it). The fix was to scale the dominating function by $t^n$
(\texttt{Integrable.const\_mul}) and finish the now-genuine equality with
\texttt{le\_of\_eq (by ring)}. A good reminder that ``mirror the $n=0$
proof'' still has to track the constants the monomial introduces.

Everything else followed the FDT-capstone template mechanically: the same
continuity-to-measurability discharges\footnote{%
\texttt{Continuous.aestronglyMeasurable} composed with \texttt{fun\_prop}.},
the same \texttt{set}-bound idiom then a \texttt{rw} closing the
integrability congruence by \texttt{rfl}.

GPT-5.5 Pro was consulted on the statement form (general-$n$ step vs.\ two
explicit derivatives, and the $(-(tx))^{n+1}$ sign). Because the proof is a
constant-factor generalisation of an already-merged result and the closed
forms were numerically confirmed against finite differences before the
build, this proceeded as a mirror-shaped tide without blocking on the
consult; the build is the correctness witness.

\section{What was Mathlib, what was new}

\textbf{Mathlib.} The dominated-differentiation
lemma\footnote{\texttt{hasDerivAt\_integral\_of\_dominated\_loc\_of\_deriv\_le}.};
\texttt{HasDerivAt.const\_mul}; the \texttt{Integrable} combinators
(\texttt{const\_mul}, \texttt{mono'}, \texttt{congr});
\texttt{integral\_congr\_ae}; and continuity-to-measurability via
\texttt{fun\_prop}.

\textbf{Seabed (reused).} The general $|x|^m$ integrability witness and the
two public pointwise helpers (pointwise derivative, pointwise bound), all
from the FDT-capstone tide; \texttt{anharmonic\_coercive}.

\textbf{New.} The general iterated-derivative step lemma and the two named
derivative corollaries. About 200 lines.

\section{Lessons}

\begin{itemize}[itemsep=3pt]
\item \textbf{Generalise the index when the supporting witness already is.}
  Because the $|x|^m$ integrability lemma is uniform in $m$, the single
  general-$n$ step lemma cost no more than $n=2$ alone and gives every
  higher derivative for free. The same observation made the FDT capstone
  cheap; it is becoming the defining move of this anharmonic arc.
\item \textbf{``Mirror the $n=0$ proof'' is not ``copy it''.} The monomial
  factor introduces a $t^n$ that the dominator must carry. The build caught
  it, but the magnitude computation should have caught it first.
\item \textbf{A constant-factor generalisation plus a numerical check is a
  legitimate proceed-without-GPT case.} The deliberation overhead is best
  spent where the statement form is genuinely uncertain, not where the proof
  is a verified template with confirmed closed forms.
\end{itemize}

\section{Follow-ups}

\begin{itemize}[itemsep=3pt]
\item \textbf{Moment normalisation.} $Z^{(n)}(0)/Z(0) = \langle(-tx)^n\rangle_0$
  ties these derivatives to the raw moments under the unperturbed Gibbs
  measure; a short corollary once divided through by $Z(0)>0$.
\item \textbf{Flow-equation / fourth-cumulant identity.} The motivating
  application: assemble $Z''(0), Z'''(0)$ (and $Z^{(4)}(0)$, a further $n=3$
  instance) into the connected fourth cumulant. Now unblocked.
\item \textbf{General-$h$ version.} Centring the dominated-differentiation
  ball at an arbitrary $h_0$ rather than $0$ would upgrade the step lemma to
  a genuine ``$Z$ is smooth in $h$'' statement and allow true iterated
  \texttt{iteratedDeriv} chaining; deferred as the $h=0$ form is what the
  cumulant application needs.
\end{itemize}

\end{document}
